% ============================================================
% Question Bank: ch01-02 Recognizing Proportional Relationships
% Topic Title: Recognizing Proportional Relationships
% CCSS: 7.RP.A.2a
% Questions: 27 (15 MC + 6 SA + 6 Graphical)
% ============================================================

% --- Multiple Choice Questions (15) ---

\begin{practiceQuestion}{1-2-q01}{mc}
\begin{questionText}
A store sells bananas. The table below shows the cost for different numbers of bananas.

\begin{center}
\begin{tabular}{|c|c|}
\hline
\textbf{Bananas} & \textbf{Cost (\$)} \\
\hline
$2$ & $0.90$ \\
$5$ & $2.25$ \\
$8$ & $3.60$ \\
\hline
\end{tabular}
\end{center}

Is this a proportional relationship?
\end{questionText}
\choiceA{Yes, because each ratio $\frac{y}{x}$ equals $0.45$}
\choiceB{No, because the cost increases by different amounts}
\choiceC{Yes, because the table has three rows}
\choiceD{No, because the ratios are not all equal}
\correctAnswer{A}
\explanation{Check each ratio: $0.90 \div 2 = 0.45$, $2.25 \div 5 = 0.45$, $3.60 \div 8 = 0.45$. Every ratio is the same, so the relationship is proportional.}
\end{practiceQuestion}

\begin{practiceQuestion}{1-2-q02}{mc}
\begin{questionText}
Which table shows a proportional relationship?
\end{questionText}
\choiceA{\begin{tabular}{|c|c|} \hline $x$ & $y$ \\ \hline $1$ & $4$ \\ $2$ & $7$ \\ $3$ & $10$ \\ \hline \end{tabular}}
\choiceB{\begin{tabular}{|c|c|} \hline $x$ & $y$ \\ \hline $2$ & $6$ \\ $4$ & $12$ \\ $6$ & $18$ \\ \hline \end{tabular}}
\choiceC{\begin{tabular}{|c|c|} \hline $x$ & $y$ \\ \hline $1$ & $5$ \\ $2$ & $8$ \\ $3$ & $11$ \\ \hline \end{tabular}}
\choiceD{\begin{tabular}{|c|c|} \hline $x$ & $y$ \\ \hline $3$ & $6$ \\ $5$ & $12$ \\ $7$ & $18$ \\ \hline \end{tabular}}
\correctAnswer{B}
\explanation{In choice B, every ratio $\frac{y}{x}$ equals $3$: $6 \div 2 = 3$, $12 \div 4 = 3$, $18 \div 6 = 3$. The other tables have ratios that change from row to row, so they are not proportional.}
\end{practiceQuestion}

\begin{practiceQuestion}{1-2-q03}{mc}
\begin{questionText}
A gym charges a $\$10$ monthly fee plus $\$5$ per visit. Which statement is true about the relationship between number of visits and total cost?
\end{questionText}
\choiceA{It is proportional because the cost increases at a constant rate of $\$5$.}
\choiceB{It is proportional because you can make a table of values.}
\choiceC{It is not proportional because the ratio of total cost to visits is not constant.}
\choiceD{It is not proportional because the cost is always more than $\$10$.}
\correctAnswer{C}
\explanation{With a $\$10$ base fee, the total cost for $1$ visit is $\$15$ ($\frac{15}{1} = 15$) and for $2$ visits is $\$20$ ($\frac{20}{2} = 10$). The ratios differ, so the relationship is not proportional.}
\end{practiceQuestion}

\begin{practiceQuestion}{1-2-q04}{mc}
\begin{questionText}
A student says the relationship in the table below is proportional because "$y$ keeps going up by $6$." What is wrong with the student's reasoning?

\begin{center}
\begin{tabular}{|c|c|}
\hline
$x$ & $y$ \\
\hline
$1$ & $8$ \\
$2$ & $14$ \\
$3$ & $20$ \\
\hline
\end{tabular}
\end{center}
\end{questionText}
\choiceA{The student confused constant difference with constant ratio.}
\choiceB{The values of $y$ do not increase by $6$.}
\choiceC{A proportional relationship always has $y$ increasing by $6$.}
\choiceD{The student should have added $x$ and $y$ instead.}
\correctAnswer{A}
\explanation{A proportional relationship requires a constant ratio $\frac{y}{x}$, not a constant difference. Here the ratios are $8$, $7$, and $6.\overline{6}$, which are not equal, so it is not proportional even though $y$ increases by $6$ each time.}
\end{practiceQuestion}

\begin{practiceQuestion}{1-2-q05}{mc}
\begin{questionText}
A graph of a relationship passes through the points $(0, 0)$ and $(4, 12)$ and forms a straight line. Which statement is true?
\end{questionText}
\choiceA{The relationship is proportional with a unit rate of $3$.}
\choiceB{The relationship is proportional with a unit rate of $4$.}
\choiceC{The relationship is not proportional because $12 \div 4 = 3$, which is odd.}
\choiceD{More information is needed to tell if it is proportional.}
\correctAnswer{A}
\explanation{A straight line through the origin means the relationship is proportional. The unit rate is $k = \frac{12}{4} = 3$.}
\end{practiceQuestion}

\begin{practiceQuestion}{1-2-q06}{mc}
\begin{questionText}
A taxi company charges $\$3$ as a base fare plus $\$2$ per mile. Which best describes the graph of total cost versus miles?
\end{questionText}
\choiceA{A straight line through $(0, 0)$}
\choiceB{A straight line through $(0, 3)$}
\choiceC{A curved line through $(0, 0)$}
\choiceD{A straight line through $(1, 2)$}
\correctAnswer{B}
\explanation{The cost equation is $c = 2m + 3$. The graph is a straight line, but it crosses the $y$-axis at $3$, not at the origin. Because it does not pass through $(0, 0)$, the relationship is not proportional.}
\end{practiceQuestion}

\begin{practiceQuestion}{1-2-q07}{mc}
\begin{questionText}
The table shows the distance a remote-control car travels over time.

\begin{center}
\begin{tabular}{|c|c|}
\hline
\textbf{Time (s)} & \textbf{Distance (ft)} \\
\hline
$3$ & $12$ \\
$6$ & $24$ \\
$9$ & $36$ \\
$12$ & $48$ \\
\hline
\end{tabular}
\end{center}

What is the best reason this relationship is proportional?
\end{questionText}
\choiceA{The distance increases by $12$ each time.}
\choiceB{The ratio $\frac{\text{distance}}{\text{time}} = 4$ for every row.}
\choiceC{The time values are all multiples of $3$.}
\choiceD{The table has four data points.}
\correctAnswer{B}
\explanation{A relationship is proportional when the ratio $\frac{y}{x}$ is the same for every pair. Here $\frac{12}{3} = \frac{24}{6} = \frac{36}{9} = \frac{48}{12} = 4$ for every row.}
\end{practiceQuestion}

\begin{practiceQuestion}{1-2-q08}{mc}
\begin{questionText}
Which of the following is \textbf{not} a requirement for a proportional relationship?
\end{questionText}
\choiceA{The graph passes through the origin.}
\choiceB{The ratio $\frac{y}{x}$ is constant.}
\choiceC{The graph is a straight line.}
\choiceD{The values of $x$ must increase by equal amounts.}
\correctAnswer{D}
\explanation{A proportional relationship requires a constant ratio and a straight line through the origin. The $x$-values do not need to increase by equal amounts --- they just need to maintain the same ratio $\frac{y}{x}$ for every pair.}
\end{practiceQuestion}

\begin{practiceQuestion}{1-2-q09}{mc}
\begin{questionText}
A bookstore sells bookmarks at a constant price with no additional fees. Which graph could represent the relationship between the number of bookmarks and the total cost?
\end{questionText}
\choiceA{A straight line starting at $(0, 2)$}
\choiceB{A curved line starting at $(0, 0)$}
\choiceC{A straight line starting at $(0, 0)$}
\choiceD{A horizontal line at $y = 5$}
\correctAnswer{C}
\explanation{Constant price with no additional fees means cost is proportional to quantity. A proportional relationship is a straight line through the origin $(0, 0)$.}
\end{practiceQuestion}

\begin{practiceQuestion}{1-2-q10}{mc}
\begin{questionText}
A pet groomer charges $\$15$ per dog with no other fees. Another groomer charges $\$12$ per dog plus a $\$6$ supply fee. Which groomer has a proportional relationship between the number of dogs and total cost?
\end{questionText}
\choiceA{Only the first groomer}
\choiceB{Only the second groomer}
\choiceC{Both groomers}
\choiceD{Neither groomer}
\correctAnswer{A}
\explanation{The first groomer's cost ($15d$) has no added fee, so the ratio is always $\$15$ per dog --- proportional. The second groomer's cost ($12d + 6$) includes a fixed fee, so the ratio $\frac{\text{cost}}{\text{dogs}}$ changes --- not proportional.}
\end{practiceQuestion}

\begin{practiceQuestion}{1-2-q11}{mc}
\begin{questionText}
Tina plots four points from a table on a coordinate plane: $(1, 3)$, $(2, 6)$, $(3, 9)$, $(4, 13)$. She says the relationship is proportional. Is she correct?
\end{questionText}
\choiceA{Yes, because the first three points line up with the origin.}
\choiceB{No, because $\frac{13}{4} \neq 3$ while the other ratios are all $3$.}
\choiceC{Yes, because most of the ratios are equal.}
\choiceD{No, because $13$ is not a multiple of $3$.}
\correctAnswer{B}
\explanation{For a proportional relationship, \textbf{every} ratio must be the same. The first three ratios are $3$, but $\frac{13}{4} = 3.25 \neq 3$. One mismatch is enough to make it not proportional.}
\end{practiceQuestion}

\begin{practiceQuestion}{1-2-q12}{mc}
\begin{questionText}
The table below represents the earnings of a student who babysits.

\begin{center}
\begin{tabular}{|c|c|}
\hline
\textbf{Hours} & \textbf{Earnings (\$)} \\
\hline
$1$ & $9$ \\
$3$ & $27$ \\
$5$ & $45$ \\
\hline
\end{tabular}
\end{center}

If the student works $0$ hours, what should the earnings be for this to be proportional?
\end{questionText}
\choiceA{$\$9$}
\choiceB{$\$0$}
\choiceC{$\$4.50$}
\choiceD{$\$1$}
\correctAnswer{B}
\explanation{For a relationship to be proportional, the graph must pass through $(0, 0)$. This means at $0$ hours worked, the earnings must be $\$0$.}
\end{practiceQuestion}

\begin{practiceQuestion}{1-2-q13}{mc}
\begin{questionText}
Carlos earns $\$8.50$ per hour washing cars with no sign-on bonus. Which equation represents his pay?
\end{questionText}
\choiceA{$p = 8.50h + 5$}
\choiceB{$p = 8.50h$}
\choiceC{$p = h + 8.50$}
\choiceD{$p = 8.50 \div h$}
\correctAnswer{B}
\explanation{With no sign-on bonus, Carlos's pay is simply the hourly rate times the hours: $p = 8.50h$. This is a proportional relationship because it has the form $y = kx$ with no added constant.}
\end{practiceQuestion}

\begin{practiceQuestion}{1-2-q14}{mc}
\begin{questionText}
A straight line on a coordinate plane passes through $(2, 6)$ and $(4, 12)$ but does \textbf{not} pass through the origin. Which could be the $y$-intercept?
\end{questionText}
\choiceA{$(0, 0)$}
\choiceB{$(0, 2)$}
\choiceC{$(0, -1)$}
\choiceD{The $y$-intercept must be $(0, 0)$ since the ratios are equal}
\correctAnswer{C}
\explanation{The slope is $\frac{12-6}{4-2} = 3$. If the line passed through $(0,0)$, at $x=2$ we'd get $y=6$, which matches. But the problem states it does not pass through the origin, so the $y$-intercept is some nonzero value such as $(0,-1)$, making it not proportional despite equal ratios in those two rows.}
\end{practiceQuestion}

\begin{practiceQuestion}{1-2-q15}{mc}
\begin{questionText}
A delivery service charges $\$4$ per package. There are no other fees. Is the total cost proportional to the number of packages?
\end{questionText}
\choiceA{Yes, because the cost per package is constant and there is no fixed fee.}
\choiceB{No, because $\$4$ is not a unit rate.}
\choiceC{Yes, but only if at least $3$ packages are shipped.}
\choiceD{No, because the cost will always be at least $\$4$.}
\correctAnswer{A}
\explanation{The total cost is $c = 4p$, which has the form $y = kx$ with no extra fee. The ratio $\frac{c}{p} = 4$ for every order, and the graph passes through $(0,0)$, so it is proportional.}
\end{practiceQuestion}

% --- Short Answer Questions (6) ---

\begin{practiceQuestion}{1-2-q16}{sa}
\begin{questionText}
A plumber charges $\$65$ per hour with no trip fee. After $4$ hours the bill is $\$260$. Is the relationship between hours and cost proportional? Write ``Yes'' or ``No.''
\end{questionText}
\correctAnswer{Yes}
\explanation{Check the ratio: $\frac{260}{4} = 65$ and the per-hour rate is $\$65$, which is constant. With no trip fee, the graph passes through $(0,0)$, confirming it is proportional.}
\end{practiceQuestion}

\begin{practiceQuestion}{1-2-q17}{sa}
\begin{questionText}
A music subscription costs $\$8$ per month plus a $\$15$ sign-up fee. Is the total cost proportional to the number of months? Write ``Yes'' or ``No.''
\end{questionText}
\correctAnswer{No}
\explanation{The $\$15$ sign-up fee is a fixed cost that does not depend on the number of months. This means at $0$ months the cost is $\$15$, not $\$0$, so the relationship is not proportional.}
\end{practiceQuestion}

\begin{practiceQuestion}{1-2-q18}{sa}
\begin{questionText}
The table shows values of $x$ and $y$.

\begin{center}
\begin{tabular}{|c|c|}
\hline
$x$ & $y$ \\
\hline
$4$ & $10$ \\
$8$ & $20$ \\
$12$ & $30$ \\
\hline
\end{tabular}
\end{center}

What is the constant ratio $\frac{y}{x}$ that shows this is proportional?
\end{questionText}
\correctAnswer{$2.5$}
\explanation{Divide $y$ by $x$ for each row: $10 \div 4 = 2.5$, $20 \div 8 = 2.5$, $30 \div 12 = 2.5$. The ratio is constant at $2.5$, confirming proportionality.}
\end{practiceQuestion}

\begin{practiceQuestion}{1-2-q19}{sa}
\begin{questionText}
Two quantities are proportional. When $x = 6$, $y = 21$. When $x = 10$, $y = 35$. What is the value of $y$ when $x = 14$?
\end{questionText}
\correctAnswer{$49$}
\explanation{Find the constant ratio: $\frac{21}{6} = 3.5$ and $\frac{35}{10} = 3.5$. Since the relationship is proportional, $y = 3.5 \times 14 = 49$.}
\end{practiceQuestion}

\begin{practiceQuestion}{1-2-q20}{sa}
\begin{questionText}
A parking garage charges $\$3$ per hour. After $5$ hours, is the total cost proportional to the number of hours? If yes, state the constant of proportionality. If no, write ``No.''
\end{questionText}
\correctAnswer{$3$}
\explanation{With a flat $\$3$ per hour and no additional fees, the total cost is $c = 3h$. The ratio $\frac{c}{h} = 3$ for every hour. The constant of proportionality is $3$.}
\end{practiceQuestion}

\begin{practiceQuestion}{1-2-q21}{sa}
\begin{questionText}
A student makes a table for a proportional relationship. Three of the entries are $(2, 7)$, $(4, 14)$, and $(6, 22)$. Is there an error? If so, which entry is wrong?
\end{questionText}
\correctAnswer{$(6, 22)$}
\explanation{Check each ratio: $\frac{7}{2} = 3.5$ and $\frac{14}{4} = 3.5$, but $\frac{22}{6} \approx 3.67 \neq 3.5$. The entry $(6, 22)$ breaks proportionality; it should be $(6, 21)$.}
\end{practiceQuestion}

% --- Graphical Multiple Choice Questions (3) ---

\begin{practiceQuestion}{1-2-q22}{gmc}
\begin{questionText}
Look at the graph below. Does it represent a proportional relationship?

\begin{center}
\begin{tikzpicture}[scale=0.65]
  \draw[step=1,gray!20,thin] (0,0) grid (7,7);
  \draw[thick,->] (0,0) -- (7.5,0) node[right]{$x$};
  \draw[thick,->] (0,0) -- (0,7.5) node[above]{$y$};
  \foreach \x in {1,...,7} \draw (\x,0.1) -- (\x,-0.1) node[below]{\tiny \x};
  \foreach \y in {1,...,7} \draw (0.1,\y) -- (-0.1,\y) node[left]{\tiny \y};
  \draw[thick,funBlue] (0,2) -- (6,6);
  \fill[funBlue] (0,2) circle (3pt);
  \fill[funBlue] (3,4) circle (3pt);
  \fill[funBlue] (6,6) circle (3pt);
\end{tikzpicture}
\end{center}
\end{questionText}
\choiceA{Yes, because it is a straight line.}
\choiceB{Yes, because the points go up at a constant rate.}
\choiceC{No, because the line does not pass through the origin.}
\choiceD{No, because the slope is less than $1$.}
\correctAnswer{C}
\explanation{A proportional relationship must be a straight line that passes through the origin $(0, 0)$. This line crosses the $y$-axis at $(0, 2)$, so it is not proportional.}
\end{practiceQuestion}

\begin{practiceQuestion}{1-2-q23}{gmc}
\begin{questionText}
The table below shows the price of lemons at a fruit stand.

\begin{center}
\begin{tabular}{|c|c|c|}
\hline
\textbf{Lemons} ($x$) & \textbf{Cost} ($y$) & $\frac{y}{x}$ \\
\hline
$3$ & $\$1.50$ & $?$ \\
$6$ & $\$3.00$ & $?$ \\
$9$ & $\$4.50$ & $?$ \\
$12$ & $\$6.00$ & $?$ \\
\hline
\end{tabular}
\end{center}

What value should fill every cell in the $\frac{y}{x}$ column?
\end{questionText}
\choiceA{$\$0.25$}
\choiceB{$\$0.50$}
\choiceC{$\$2.00$}
\choiceD{$\$1.50$}
\correctAnswer{B}
\explanation{Divide cost by lemons for any row: $1.50 \div 3 = 0.50$, $3.00 \div 6 = 0.50$, $4.50 \div 9 = 0.50$, $6.00 \div 12 = 0.50$. The constant ratio is $\$0.50$ per lemon.}
\end{practiceQuestion}

\begin{practiceQuestion}{1-2-q24}{gmc}
\begin{questionText}
Two lines are graphed on the same coordinate plane. Line $A$ passes through $(0, 0)$ and $(5, 15)$. Line $B$ passes through $(0, 4)$ and $(5, 19)$.

\begin{center}
\begin{tikzpicture}[scale=0.32]
  \draw[step=2,gray!20,thin] (0,0) grid (10,20);
  \draw[thick,->] (0,0) -- (10.5,0) node[right]{$x$};
  \draw[thick,->] (0,0) -- (0,21) node[above]{$y$};
  \foreach \x in {2,4,6,8,10} \draw (\x,0.2) -- (\x,-0.2) node[below]{\tiny \x};
  \foreach \y in {4,8,12,16,20} \draw (0.2,\y) -- (-0.2,\y) node[left]{\tiny \y};
  \draw[thick,funBlue] (0,0) -- (7,21);
  \draw[thick,funRed] (0,4) -- (6,22);
  \node[right,funBlue] at (6,18) {\small $A$};
  \node[right,funRed] at (5.5,20.5) {\small $B$};
\end{tikzpicture}
\end{center}

Which line represents a proportional relationship?
\end{questionText}
\choiceA{Only Line $A$}
\choiceB{Only Line $B$}
\choiceC{Both lines}
\choiceD{Neither line}
\correctAnswer{A}
\explanation{Line $A$ passes through the origin $(0, 0)$ and is straight, so it is proportional. Line $B$ passes through $(0, 4)$, which means it has a $y$-intercept of $4$ and is not proportional.}
\end{practiceQuestion}

% --- Graphical Short Answer Questions (3) ---

\begin{practiceQuestion}{1-2-q25}{gsa}
\begin{questionText}
The table below shows the relationship between the number of tickets and the total cost at a carnival.

\begin{center}
\begin{tabular}{|c|c|}
\hline
\textbf{Tickets} & \textbf{Total Cost (\$)} \\
\hline
$0$ & $0$ \\
$3$ & $7.50$ \\
$5$ & $12.50$ \\
$8$ & $20.00$ \\
\hline
\end{tabular}
\end{center}

Is the relationship proportional? If yes, state the unit rate. If no, write ``No.''
\end{questionText}
\correctAnswer{$\$2.50$}
\explanation{Check: $7.50 \div 3 = 2.50$, $12.50 \div 5 = 2.50$, $20.00 \div 8 = 2.50$. All ratios equal $2.50$, and the table includes $(0, 0)$. The relationship is proportional with a unit rate of $\$2.50$ per ticket.}
\end{practiceQuestion}

\begin{practiceQuestion}{1-2-q26}{gsa}
\begin{questionText}
The scatter plot below shows the heights of plants (in cm) after different numbers of weeks.

\begin{center}
\begin{tikzpicture}[scale=0.55]
  \draw[step=1,gray!20,thin] (0,0) grid (8,10);
  \draw[thick,->] (0,0) -- (8.5,0) node[right]{\small Weeks};
  \draw[thick,->] (0,0) -- (0,10.5) node[above]{\small Height (cm)};
  \foreach \x in {1,...,8} \draw (\x,0.1) -- (\x,-0.1) node[below]{\tiny \x};
  \foreach \y in {1,...,10} \draw (0.1,\y) -- (-0.1,\y) node[left]{\tiny \y};
  \fill[funGreen] (0,0) circle (4pt);
  \fill[funGreen] (2,3) circle (4pt);
  \fill[funGreen] (4,6) circle (4pt);
  \fill[funGreen] (6,9) circle (4pt);
\end{tikzpicture}
\end{center}

Is this a proportional relationship? If yes, what is the growth rate per week?
\end{questionText}
\correctAnswer{$1.5$ cm per week}
\explanation{The points $(0,0)$, $(2,3)$, $(4,6)$, and $(6,9)$ lie on a straight line through the origin. The ratio $\frac{y}{x}$ is $\frac{3}{2} = 1.5$ for every point. The relationship is proportional at $1.5$ cm per week.}
\end{practiceQuestion}

\begin{practiceQuestion}{1-2-q27}{gsa}
\begin{questionText}
The graph below shows the cost of renting kayaks.

\begin{center}
\begin{tikzpicture}[scale=0.5]
  \draw[step=1,gray!20,thin] (0,0) grid (8,10);
  \draw[thick,->] (0,0) -- (8.5,0) node[right]{\small Hours};
  \draw[thick,->] (0,0) -- (0,10.5) node[above]{\small Cost (\$)};
  \foreach \x in {1,...,8} \draw (\x,0.1) -- (\x,-0.1) node[below]{\tiny \x};
  \foreach \y in {2,4,...,10} \draw (0.1,\y) -- (-0.1,\y) node[left]{\tiny \y0};
  \fill[funPurple] (0,2) circle (4pt);
  \fill[funPurple] (1,3.5) circle (4pt);
  \fill[funPurple] (2,5) circle (4pt);
  \fill[funPurple] (4,8) circle (4pt);
  \draw[thick,funPurple] (0,2) -- (4,8);
\end{tikzpicture}
\end{center}

The $y$-axis values are in tens of dollars ($20, 40, \ldots$). Why is this relationship \textbf{not} proportional? State the reason.
\end{questionText}
\correctAnswer{The line does not pass through the origin}
\explanation{The graph starts at $(0, 20)$, meaning there is a $\$20$ base fee even for $0$ hours. A proportional relationship must pass through $(0, 0)$. Since the line starts above the origin, the relationship is not proportional.}
\end{practiceQuestion}
